\section{Independent channels: exact statement and limited consequence}
\label{sec:channel_theorem}
\label{sec:channel-theorem}

We now isolate the part of the analysis that is a theorem.  Let
\begin{equation}
 \mathcal V=\operatorname{Hom}(P\mathcal H,Q\mathcal H)
 \label{eq:response_space}
\end{equation}
at a fixed base point.  Cumulants of a complex matrix in $\mathcal V$ are
defined on the underlying real vector space; complex trace words are obtained
by complexifying the resulting real multilinear tensors.  This convention
retains holomorphic, antiholomorphic, and mixed contractions and is
equivalent to retaining both tensors in
Eq.~\eqref{eq:covariance_pseudocovariance}.

\paragraph{Theorem 1 (weighted cumulant additivity).}
Let $Y_1,\ldots,Y_n\in\mathcal V$ be mutually independent centered complex
random matrices, and let $w_1,\ldots,w_n\in\mathbb R$ be deterministic.  For
an order-$k$ statement, assume that every real-linear projection under
consideration has a finite $k$th absolute moment.  All channels must occupy
the same response space, or be identified with it by a fixed map independent
of the outcome.  Then, for $Z$ in
Eq.~\eqref{eq:random_channel_model},
\begin{equation}
 \boxed{
 \kappa_k(Z)=\sum_{\alpha=1}^{n}
 w_\alpha^k\kappa_k(Y_\alpha)} .
 \label{eq:cumulant_additivity}
\end{equation}
The equality is a finite-$n$ tensor identity.  It is not a central-limit
approximation.

To see the result, let $t$ be any real-linear test field on $\mathcal V$.
Independence factorizes the characteristic function and hence its logarithm,
\begin{equation}
 \log\mathbb E e^{it(Z)}
 =\sum_{\alpha=1}^{n}
 \log\mathbb E e^{iw_\alpha t(Y_\alpha)} .
 \label{eq:log_characteristic_additivity}
\end{equation}
Taking $k$ derivatives at the origin supplies $w_\alpha^k$ and produces no
mixed-channel connected term.  Polarization over distinct test fields proves
Eq.~\eqref{eq:cumulant_additivity}.  A detailed realification proof is given
in the Supplemental Material.

At second order, with $y_\alpha=\operatorname{vec}(Y_\alpha)$, the ordinary
covariance and pseudocovariance are
\begin{equation}
 C_\alpha=\mathbb E[y_\alpha y_\alpha^\dagger],
 \qquad
 \Pi_\alpha=\mathbb E[y_\alpha y_\alpha^T],
 \label{eq:channel_second_moments}
\end{equation}
and independence gives
\begin{equation}
 C_Z=\sum_\alpha w_\alpha^2C_\alpha,
 \qquad
 \Pi_Z=\sum_\alpha w_\alpha^2\Pi_\alpha.
 \label{eq:weighted_second_moments}
\end{equation}
The connected fourth tensor must subtract the Wick tensor constructed from
both $C_Z$ and $\Pi_Z$.  Pairwise uncorrelated channels are insufficient for
Theorem~1: mixed fourth cumulants can remain when independence fails.

The participation number of the deterministic weights is
\begin{equation}
 N_{\mathrm{eff}}
 =\frac{\left(\sum_\alpha w_\alpha^2\right)^2}
 {\sum_\alpha w_\alpha^4} .
 \label{eq:effective_channel_number}
\end{equation}
For $L^2$-normalized equal weights it equals $n$.  Concentrated weights can
make it much smaller than the nominal channel count.  Theorem~1 alone does
not produce an inverse-$N_{\mathrm{eff}}$ law because the tensors
$\kappa_4(Y_\alpha)$ need not be equal or even comparable.

\paragraph{Corollary 1 (conditional fourth-cumulant law).}
Assume in addition that the tested channel contraction is variance
normalized, every channel has a finite fourth moment, and the standardized
fourth cumulants are identical,
$\kappa_4(Y_\alpha)=K_4$.  Then
\begin{equation}
 \frac{\kappa_4(Z)}
 {\left(\sum_\alpha w_\alpha^2\right)^2}
 =\frac{K_4}{N_{\mathrm{eff}}} .
 \label{eq:inverse_neff_law}
\end{equation}
If instead $\|\kappa_4(Y_\alpha)\|\le K$ in one fixed tensor norm, then
\begin{equation}
 \frac{\|\kappa_4(Z)\|}
 {\left(\sum_\alpha w_\alpha^2\right)^2}
 \le\frac{K}{N_{\mathrm{eff}}} .
 \label{eq:inverse_neff_bound}
\end{equation}
For unequal known channel cumulants, the executable prediction is
\begin{equation}
 \kappa_{4,\mathrm{pred}}(Z)
 =\frac{\sum_\alpha w_\alpha^4\kappa_4(Y_\alpha)}
 {\left(\sum_\alpha w_\alpha^2\right)^2} .
 \label{eq:unequal_channel_prediction}
\end{equation}
No exponent is fitted in Eqs.~\eqref{eq:inverse_neff_law}--
\eqref{eq:unequal_channel_prediction}.

The analytic implementation was checked with a fixed non-Gaussian complex
matrix population having zero mean, unit variance, zero population
pseudocovariance, finite fourth moment, and standardized fourth cumulant
$K_4=8$.  Equal and linearly unequal weights were evaluated at
$n=4,8,16,32,64$ using $50{,}000$ matrices and $25$ uncertainty blocks.  The
largest direct relative error from
Eq.~\eqref{eq:unequal_channel_prediction} was $\ChannelTheoryMaximumRelativeError$, below the
registered $0.08$ implementation threshold.  This oracle verifies the code
path under data generated from the theorem assumptions.  It is not evidence
that the response channels of a physical Hamiltonian are independent.

The second exact statement concerns only accessibility of the fiber.
Given response matrices $X_a$, define the unital star-closed algebra
\begin{equation}
 \mathcal A_X
 =\operatorname{alg}^{*}\{X_a^\dagger X_b:a,b\}
 \subseteq\operatorname{End}_{\mathbb C}(P\mathcal H)
 \label{eq:response_algebra}
\end{equation}
and its commutant
\begin{equation}
 \mathcal A_X'
 =\{M:[M,A]=0\ \text{for every }A\in\mathcal A_X\} .
 \label{eq:response_commutant}
\end{equation}

\paragraph{Proposition 1 (commutant criterion).}
For a finite-dimensional complex star algebra,
$\mathcal A_X'=\mathbb C I_D$ implies that $\mathcal A_X$ acts irreducibly on
the protected fiber.  Burnside's theorem then gives
\begin{equation}
 \mathcal A_X=\operatorname{End}_{\mathbb C}(P\mathcal H).
 \label{eq:burnside_response}
\end{equation}
A nonscalar commutant therefore rules out an unrestricted full-matrix mixing
ansatz until its invariant blocks are resolved.  The converse use relevant
to statistics is forbidden: a scalar commutant does not imply independent
channels, Gaussianity, suppression of connected cumulants, locality,
thermalization, or ETH.  A deterministic pair of matrices can generate the
full matrix algebra.

The theorem and proposition leave several explicit failure branches.  Mixed
connected tensors modify Eq.~\eqref{eq:cumulant_additivity} for correlated
channels.  Without finite fourth moments, the statistic in
Eq.~\eqref{eq:inverse_neff_law} is undefined.  If
$\max_\alpha w_\alpha^2/\sum_\beta w_\beta^2$ does not decrease, nominal
channel proliferation does not imply growing $N_{\mathrm{eff}}$.  Unresolved
symmetry blocks produce a nonscalar commutant; nonstationary tangent
populations invalidate a comparison based only on $N_{\mathrm{eff}}$; and a
closing external gap invalidates the isolated-projector response.  Finally,
omitting pseudocovariance removes the third Wick pairing and changes the null
being tested.

These limitations are active in the result boundary.  The analytic audit and
exact index construction passed, but the prospectively sealed chiral random
tangent validation did not remain inside its fixed channel-law band.  The
selected branch is therefore \texttt{random\_channel\_failure}.  Retrospective
data from other protected models lack documented independent response units
and cannot repair that prospective failure.  We retain
Eqs.~\eqref{eq:cumulant_additivity}--\eqref{eq:burnside_response} as exact
conditional statements and reject the stronger inference that they establish
Geometric ETH for exactly degenerate quantum-state bundles.
